% 中文审阅稿 —— 由 paper/build_paper_zh.py 生成，请勿手改。
% 英文版 paper/main.tex 是主张的唯一来源；本文件是派生出的审阅副本。
\documentclass[11pt]{ctexart}
\usepackage{booktabs}
\usepackage{amsmath}
\usepackage{amssymb}
\usepackage[margin=2.4cm]{geometry}
\usepackage[hidelinks]{hyperref}

\newcommand{\zhbanner}[1]{\noindent\fbox{\parbox{\dimexpr\linewidth-2\fboxsep-2\fboxrule}{#1}}}
\title{\bf 面向 LLM 智能体的旁路参数化记忆：\\可精确擦除的槽位，以及不占上下文的读取}
\author{成昊天 \\ \small 南京航空航天大学}
\date{2026 年 9 月 30 日}

\begin{document}
\maketitle

\zhbanner{\small\textbf{这是中文审阅稿。}它与英文版 \texttt{main.pdf} 逐节对应，
所有数字都由 \texttt{runs/summary.json} 生成，与英文版同一来源；
\textbf{技术主张以英文版为准}。请先看三处：（一）§\ref{sec:neg} 里 B1 保留项
“未复现”的措辞；（二）§\ref{sec:limits} 关于“质量优势随规模收窄”的表述；
（三）表格是否一次就能读懂。}

\vspace{0.6em}
\begin{center}\small
数字来源：\texttt{runs/summary.json}。生成脚本：\texttt{paper/build\_paper\_zh.py}。
\end{center}

\section{摘要}
LLM 智能体的长期记忆放在上下文窗口里，既要挤占推理所需的 token，又与当前工作集无法区分。
我们换了一条路：在冻结主干上挂一组容量受限的低秩记忆槽，用梯度写入、经前向传播读取、
靠擦除一个秩块遗忘。在 0.6B 模型上，写入只落到自己的槽；被擦除的槽与从未写过的槽逐位相同，
30/30 次驱逐与一次进程重启都如此；真正的约束是读时组合：把槽全加起来召回为
0/40，按查询键相似度只选一个则恢复到 35/40（oracle 38/40），两者之差恰好等于
路由错误。无标签的自我核对用 5 次写入达到全写的召回（全写 18 次），
基座漂移 0.54 nats（全写 2.59 nats）。有两个假设失败了，我们如实报告：
基于似然阈值的惊讶度判据量到的是措辞，参数化记忆也没有减少上下文惯性。在 30
道 LongMemEval single-session 题（历史完全不在场）上，参数臂复现 1.000 的参考答案
且不额外占用 prompt token，而把同样信息放进上下文只能答对 0.867 到
0.967，并付出 117 到 1688 个 token。
另外两项对照值得一提：在\textbf{条目-梯度步数相同}的条件下，槽位方案能复现这 30 条记忆中的
1.000，而全量微调只到 0.133、等秩单适配器只到 0.000；
而\textbf{只用模型之外的手段}换一个检索键，就能把 1.7B 的路由从 0.719 抬到
1.000。

\section{引言}

\subsection{上下文通道的两个代价}
在不改动权重的前提下，新信息只能进上下文窗口。这条通道有三个结构性限制。
它是零和的：记住的 token 不再是可用来推理的 token。它不可区分：每句被记住的话
都同等在场，长期记忆与当前工作集没有分别。它也退不出去：内容一旦进入上下文，
除非重写提示词并丢掉 KV cache，否则无法有选择地移除。

\subsection{一个后来被我们否证的动机}
我们当初的假设是：上下文里的内容会被模型当作当前正在发生的事。我们想知道，活在权重里的记忆是否更不容易
被带过话题边界。这个假设写在 §\ref{sec:neg}，
因为我们在那里否证了它。对一个读者来说，这个否证比一个含糊的正面结论更有用。

\subsection{我们造了什么}
一个冻结主干，加一条 $K$ 个槽的旁路。每个投影层挂一个 LoRA 适配器，总秩
$R = K \cdot r$ 被切成 $K$ 个连续块，第 $k$ 块就是第 $k$ 个槽。写入可以只作用于一个块，
读取可以只加回一个块，擦除就是把该块的初始 $A_k$ 恢复，把 $B_k$ 清零。
擦除因此成了一个可以逐位验证的命题，而不再是一个大致的感觉。

\subsection{贡献}
\begin{enumerate}\setlength\itemsep{2pt}
\item 一条秩块旁路，三种操作都能独立检查：写入被断言只触碰一个槽，擦除被断言把槽还原成
从未写过的逐位状态，读取可以在逐字节相同的模型上逐槽消融。这三条在容量上限处与进程重启后同样成立。
\item 读时组合是约束所在：全部相加毫无价值，只选一个槽约等于 oracle，选两个就损失大部分收益；
而 oracle 与 top-1 之差恰好等于路由错误，在我们测的两个规模上都如此。
\item 一个无标签的写入判据，用 5 次写入达到全写的召回（全写 18 次），
基座漂移 0.54 nats（全写 2.59 nats），并给出它所替代的似然判据的否证。
\item 一个范围受限的真实内容诊断：同一个路由器把同样的信息送进上下文，由与人工标注校准过的
LLM 判官打分；参数臂在零额外 token 下全部答对，而上下文递送要付
117 到 1688 个 token。
\item 我们报告而没有埋掉的负结果：参数化记忆没有减轻上下文惯性；似然阈值量到的是措辞；
以及一个在 $n = 24$ 时看起来修好了两槽召回的保留项，
在 $n = 64$ 时没有复现，尽管它确实把基座损伤从
1.518 降到 0.317 nats。
\end{enumerate}

\section{相关工作}
\label{sec:related}

\paragraph{运行时参数化记忆。}Titans 是最近的先例：一个在测试时按“惊讶度”更新的神经长期记忆，
其衰减起到遗忘作用。它改动主干，我们冻结主干。GradMem 用测试时梯度把上下文信息写进
memory token，是与我们设定最接近的工作，但它没有遗忘机制，而那正是我们测得最仔细的一维。
TTT 一类方法把上下文压缩进权重，并报告长上下文下的延迟收益；其中一篇给出一个数：
在 128K 上，全注意力的朴素逐 token 开销是它所需开销的 $2.7$ 倍。这些工作都没有提供逐条擦除。

\paragraph{槽位与适配器。}ProCL 把 LoRA 适配器组织成结构化的“程序记忆”槽，按输入条件注意力检索，
动机来自互补学习系统，在架构上是与我们槽位想法最近的邻居。它的设定是有任务数据集的任务增量微调，
不遗忘；它的槽是彼此独立的适配器，而不是同一个适配器内部的秩块。WISE 观察到终身编辑的
“不可能三角”（可靠性、泛化性、局部性），并用主副记忆分离来应对，那正是我们实例化的结构。
AlphaEdit 把编辑投影到被保护知识的零空间，我们在锚定项里逐槽使用了这个技术。
它们都没有提供一种能在运行时、在冻结主干上被独立写入、读取和擦除的槽；
我们测的就是这个组合，而其中的擦除一半，正是“逐位相同”这类主张得以成立的原因。

\paragraph{事实编辑。}ROME 与 MEMIT 一类编辑器离线把事实写进权重，评价指标主要是涟漪一致性；
有工作报告说一个简单的上下文编辑基线就能在这个指标上超过它们，这个警告我们认真对待
（见 §\ref{sec:limits}）。因此我们的协议把这条基线明确带上：§\ref{sec:results} 里的上下文臂
就是用编辑提示词的方式把同样的事实交给模型，而它们正是我们报告的对照。

\paragraph{智能体记忆。}MemGPT/Letta、Mem0、A-MEM、MemoryBank、MemoRAG 与 Zep 是实践中
真正被使用的系统。它们都把记忆存在上下文层或外部存储里，再检索回提示词；
没有一个做运行时参数更新。我们的设定因此是空着的，这也意味着没有现成的基线协议可以继承。
§\ref{sec:protocol} 定义了一套。

\section{方法}
\label{sec:method}

\subsection{秩块槽位与精确擦除}
对每个目标投影，$\Delta W = \frac{\alpha}{r} B A$，其中 $A$ 按秩块划分。
槽 $k$ 只能写 $A$ 的第 $k$ 块与 $B$ 的第 $k$ 列，读掩码决定加回哪些块。
擦除就是恢复 $A_k^{\text{init}}$ 并把 $B_k$ 置零，所以擦除与从未写过
是可以逐位比较的两件事。

\subsection{写入目标}
写入损失有三项：
\begin{equation}
\mathcal{L} = \mathrm{CE}(v \mid q)
\; + \; \lambda\, D_{\mathrm{KL}}\!\left(p_{\text{frozen}} \,\|\, p_{\text{slot}}\right)
\; + \; \lambda_{\mathrm{ret}}\, D_{\mathrm{KL}}\!\left(p_{\text{bank}} \,\|\, p_{\text{bank}+k}\right).
\end{equation}
第一项只开本槽，学的是查询到答案这条关联。第二项在一组通用问题上把分布拉回冻结态，
两处细节都是被测量逼出来的：它必须\textbf{逐槽}评估 —— 按已写槽求和评估时各槽位移会互相抵消，
4 个槽在那种配置下求和 $D_{\mathrm{KL}}$ 只有 0.080，
而逐个读却合计 1.790（相差 22.3x），
于是这一项在"集体"意义上被满足、每个槽单独看却毫无约束；
\textbf{改成逐槽之后}同一测量是 0.387 对 0.469（1.2x），
剩下的才是诚实的残留。它还必须覆盖每个位置：只看末位时单步散度会被压得极低，
自回归解码仍然复合成复读机。第三项是保留项（B1）。参考分布取自写入之前的记忆库，
那时已写槽激活、新槽零初始化因而还不贡献；查询用的是更早那些记忆的查询，
当前项把新槽加上去，也就是干扰真正发生的那个读条件。它的效果见 §\ref{sec:neg}：
方向对了，但没有复现。

\subsection{读时选择}
所有槽一起激活时答案会崩。原因不难理解：每个槽都被训练成产生一个低熵答案，
它们的和是一个被支配的混合。各槽本身完全可分，读错的槽什么都得不到，
所以读取路径应当选择。我们用查询键（查询在冻结主干下的末位隐状态，取键时关闭记忆，
避免键随已写记忆漂移）与各槽键做余弦相似度，只激活最匹配的那一个。这一步不新增参数，
也不需要训练。

\subsection{什么值得写}
四种判据：\texttt{always} 全写；\texttt{surprise} 看写入前损失是否超阈值；
\texttt{selfcheck} 先问冻结模型是否知道这个内容；\texttt{explicit} 只写带标记的内容。

\subsection{遗忘与持久化}
容量 $K$ 加上 5 种淘汰策略（FIFO、LRU、LFU、效用、效用加时间衰减），
另有重要标记作为硬保护。持久化用 safetensors 快照加元数据严格校验：
槽数、秩、alpha、模块集合与形状必须一致，不一致就报错，不做静默加载。

\section{评测协议}
\label{sec:protocol}

\paragraph{虚构基准。}每个实体与取值都由音节生成，因此答案不可能来自预训练知识；
否则答对无法归因于写入。同理，已经知道这一类要先问过模型：
13 条真实常识里，模型把 $2+2$ 答成 2，把一周的天数答成 1，
把最大的行星答成 Mars，最后只有 13 条可用。

\paragraph{五段隔离检查（P1 到 P5）。}P1 驱逐：答案不得出现在读时上下文里
（精确、规范化、数值三种匹配），否则该条不计入。P2 反事实：取值按种子随机化。
P3 负对照：从未写过的实体不得变好。P4 消融：上下文逐字节相同，只改活跃槽集合。
P5 仅提示词：剥掉提示词格式带来的收益。只有在配对差异的 bootstrap 置信区间不含零时，
我们才报告主张。

\paragraph{改写探针。}读取用不同措辞的问法，记忆仍按原始措辞写入。
若沿用同一个字符串，路由准确率会到 100\%，那测到的只是字符串匹配。

\paragraph{地板筛选。}惯性实验先测没有任何前提的条件，把模型本来就能答对的场景剔除。
25/30 个场景存活，被剔除的 5 个都是路边情形。

\paragraph{真实内容，以及怎么判分。}公开基准诊断用 LongMemEval 的 oracle 划分。
每条记忆就是基准自带的问题与答案对，证据轮由基准自己的 \texttt{has\_answer} 标记定位，
检索因此被 oracle 化。对照组拿到同样的信息，由同一个查询键路由器选出来放进上下文，
唯一的变量是介质。判分走语义。词面包含判据奖励逐字复现参考答案，
而这正是训练过的写入擅长的事，它无法比较改写作答的臂。所以真实内容结果由一个托管的
LLM 判官打分，一次只判一个候选，并与 10/10 条人工标注校准。
标注样本匿名且打乱，判官的结论在人工标完之后才展示。我们同时报告校准规模与
判官提示词改过一次这两件事，因为它们限定了这个一致率值多少。另有两点：
同一批题上的比较一律用配对检验，也就是精确符号检验；生成预算被当作处理的一部分，
所有臂用同一个上限，且这个上限要能让最长的参考答案写得下。

\section{结果}
\label{sec:results}

\subsection{写入只落到该落的槽（T1）}
写入隔离断言处处通过，被擦除的条件与仅提示词的条件逐字符相同（10/10），
写完的目标交叉熵是 0.0000。表里那个很低的 \texttt{mem\_on} 行是求和式读取的结果，
它属于 §\ref{sec:results} 要处理的问题，不能当作写入失败的证据。

\subsection{读时组合压倒一切（T2）}
\begin{table}[h]\centering\small
\begin{tabular}{lccc}\toprule
臂 & 0.6B & \textbf{1.7B} \\ \midrule
\texttt{oracle}（只开本槽） & 38/40 & 0.950 \\
\texttt{top1}（路由最优槽） & 35/40 & 0.675 \\
\texttt{top2}（路由前两槽） & 15/40 & 0.400 \\
\texttt{all}（全部相加） & 0/40 & 0.000 \\
路由准确率 & 0.925 & 0.725 \\ \bottomrule
\end{tabular}
\caption{改写问法下的读时组合，每种设置 40 个探针。}
\end{table}

写入同样的记忆、只改哪些槽参与：全部相加是 0/40，只选一个槽是 35/40
（oracle 38/40），选两个已经掉到 15/40。算术在这里精确闭合，
oracle 减 top-1 恰好等于路由错误，说明选择规则本身不引入损失。实用结论比
越稀疏越好更强：这种记忆的读取路径应当只选一个槽，top-$k$ 的对冲会损失大部分收益。

这个形状在 1.7B 上同样出现：全部相加仍是 0.000，只选一个槽是
0.675（oracle 0.950），选两个是 0.400，
同样的算术依然闭合，oracle 减 top-1 等于 $1 - 0.725$。
随规模改变的是损失出现的位置。路由器在 1.7B 上明显更差，0.6B 是 0.925，
1.7B 是 0.725，与 oracle 的差距于是主要来自检索。
下一步该改进的部件是路由器。

\subsection{限制这道差距的是键，而不是槽位机制}
既然路由准确率一对一变成绩效，那就值得先问一个不用训练任何参数的问题：\textbf{现行的键是不是选错了？}
答案是选错了。同一批写入、同一批探针，\textbf{只改检索键的定义}，top1 就从 0.938 抬到
1.000（0.6B）、从 0.719 抬到 1.000（1.7B），
端到端召回同步上升。胜出的两种定义都是\textbf{保留更多查询信息}的做法：最佳的是 lexical
（查询与已存查询之间的词面 TF-IDF 重叠），把它与隐状态余弦五五混合紧随其后。
有两种做法输给了现行键，都值得写出来：\textbf{浅层的键}丢掉了大部分信号（0.6B 上只有
0.312）；\textbf{用模型自己的置信度去重排前两名候选}比直接相信键\textbf{更差}
（0.688）—— 一个槽"对自己更有把握"并不构成"它才是该回答的那个槽"的证据。
按同一批探针做配对检验，这个提升在 1.7B 上显著（9 帮助、
0 损害，$p$ 为 0.0039），在 0.6B 上不显著
（2 对 0，$p$ 为 0.5000）：
\textbf{模型越大，键越重要}。

\textbf{接下来读者会想到的两条路 —— 训练一个键、以及把多个键投票聚合 —— 都没有换来应有的收益。}
用一个自监督的投影（对规范问法做增强视图、用 InfoNCE 训练，\textbf{改写过的探针问法在训练中从未出现}），
在更大的基座上 top1 是 0.906、端到端 0.844，
而它本要改进的原始隐状态是 0.719、它本要超越的纯词面信号是 1.000；
在更小的基座上它\textbf{干脆比原始键还差}（0.812 对 0.938），
配对记录也是负的：帮到 0 条、损害 3 条
（$p$ 为 0.2500）。
也就是说，\textbf{当基座大到足以让原始键变弱时，训练确实带来了真实改进，却仍然够不到那个不需要训练的信号}。
把四个键做 Borda 投票聚合，在两个尺度上都只是\textbf{恰好追平最好的单键}
（1.000 与 1.000），没有增量。
这与本节其余部分是同一个模式：\textbf{活下来的是最朴素的那个信号。}

这里有一条必须说明的边界：这些查询来自我们的\textbf{合成生成器}，内容词高度区分；
而真实基准里同一实体的多个问题\textbf{共享几乎全部词面}，所以"词面成分在真实内容上是否有用"
是一个独立的问题，见 §\ref{sec:limits}。

\subsection{修好路由之后：路由差距消失了，碰撞还在}
把组合实验\textbf{换成模型之外的键重跑}，就把上面那笔算术从"相关"变成了\textbf{因果检验} —— 而且\textbf{两个尺度上都成立}。
oracle 几乎不受换键影响（0.6B 是 0.950、1.7B 是 0.950），
单槽答对 0.950 与 0.950，路由都达到 1.000，
而 \textbf{oracle 减 top1 的差距降到 0.000（0.6B）与 0.000（1.7B）} ——
在旧键下它们分别是 $1-0.925$ 与 $1-0.725$。
也就是说，那道差距\textbf{全部来自检索}，与槽位无关。
反过来，\textbf{两槽的代价完全不受这次修理影响}：0.6B 是 0.500 对 0.950，
1.7B 是 0.525 对 0.950，而全开在两种键下都只值 0.000。
这是本节结论最锋利的形态：\textbf{换个更好的键能修的是路由，修不了的是"同时读两条记忆"}。

\subsection{真实内容上同样会崩}
合成数据的查询词高度区分、取值只有一句话，所以自然要问：\textbf{崩塌是机制的性质，还是生成器的性质？}
我们把 60 道 LongMemEval single-session 题\textbf{一题一槽}写进去（与真实内容诊断完全一样），
然后让\textbf{同时打开的槽数}逐步增加来读回（见 Table~\ref{tab:t7collision}）。
检索全程 oracle，所以唯一变化的是\textbf{读取路径同时持有多少条记忆}：
包含率从"只开正确那一槽"的 1.000，降到多开一个的 0.883、
两个的 0.733、四个的 0.483，再到整库全开的 0.000
（冻结地板是 0.017）。
⇒ \textbf{机制在真实内容上的表现与在生成内容上一致}。有一个差别要如实说：
\textbf{两槽那一步在真实内容上更温和}，也就是说\textbf{生成器放大了效应的大小，却没有凭空制造它}。
每个剂量上效应都是\textbf{单向的}（多开槽\textbf{从未}改善任何一题）：配对计数在"多开一个"这一步是
7 损害对 0 帮助（$p$ 为 0.0156），
"多开两个"是 16 对 0（$p$ 为 0.0000），
整库全开是 60 对 0（$p$ 为 0.0000）。
\textbf{两槽这一步是整条曲线上最小的效应，而它在 60 题上已经显著} ——
\textbf{"只读了几道题"这个解释因此被排除}。

\subsection{该写什么（T3）}
把 $5$ 条虚构事实与 13 条已验证常识混成一条流。全写需要
18 次梯度写入，其中 13 次花在模型本来就会的事实上，
基座漂移 2.59 nats。自我核对只用 5 次写入，浪费 0，
漏掉 0 个目标，漂移 0.54 nats。召回相同，写入少 72\%，
漂移少 76\% 到 84\%。后一个范围来自 3 个种子，
因为写入次数是确定的，漂移随种子变化。这条判据在规模控制下\textbf{连计数都没变}：
1.7B 上同样的对比是 5 次写入对 22、零漏写，
漂移 0.45 对 2.97，即 6.6x 倍，
而 0.6B 上是 4.8x 倍 —— 基座越大，这道保险\textbf{越值钱}。
似然判据是反面教材，它几乎分不开两类
（见 §\ref{sec:neg}），因为固定措辞量到的是模型会不会用这个说法，与它是否知道这件事无关。

\subsection{遗忘是精确的（T4）}
12 条记忆流进 6 个槽，前 $3$ 个被反复访问。
每次驱逐都把槽还原为初始态：30/30 个槽被验证与从未写过逐位相同，
每种策略各 6/6 个，且\textbf{在 3 个种子上完全一致} ——
那些是逐张量断言，本来就确定。被标记为重要的槽扛过了每一次溢出。FIFO 只留下
1/1 个正在被使用的记忆，它专挑你在用的丢、留着你没碰过的；
其余四种策略都留下 3/3，这个差距在每个种子上都一样（100\%--100\%）。
\textbf{真正不稳的是我们原先对那四种策略的说法}：它们在"正在被使用"的记忆上打平，
但在\textbf{没被使用}的记忆上，三种里有三种在三个种子中的一个种子上掉了一个存活者，
而 \textbf{LFU 每个种子都全保}。这既是对"简单 LFU 难以被超越"的小小独立复现，
也提醒一句：\textbf{只观察一次的"打平"未必是打平。}
在更大的基座上这五种策略表现得一模一样，\textbf{计数分毫不动}：FIFO 依然只留下 1/1 个
正在被使用的记忆，保序策略留下 3/3，而 30/30 个槽被逐位验证为从未写过。

\subsection{跨进程（T6）}
新进程里，加载前 \texttt{virgin} 为 \texttt{True}；重启后召回 6/6，
路由召回 5/6，与路由准确率 5/6 相同，所以差的那一题属于路由错误；
加载后擦除仍逐位精确（\texttt{True}）。8 个槽的快照占
18.4\,MB。

\subsection{真实内容：为什么不把记忆放进上下文（T7）}
30 道 LongMemEval single-session 题，历史完全不在场，同一批题、同一个路由器、语义判官打分：

\begin{table}[h]\centering\small
\begin{tabular}{lccc}\toprule
臂 & 判官正确率（0.6B） & 判官正确率（1.7B） & 额外 prompt token \\ \midrule
\textbf{参数记忆} & \textbf{1.000} & \textbf{1.000} & \textbf{0} \\
30 条全塞上下文 & 0.967 & 0.967 & 1688 \\
检索 1 条 / 3 条 & 0.900 / 0.900 & --- & 117 / 219 \\
oracle 单条 & 0.867 & 0.967 & 117 \\
冻结（地板） & 0.033 & 0.100 & 60 \\ \bottomrule
\end{tabular}
\caption{真实内容上的对比，判官已与人工标注校准 10/10。}
\end{table}

参数臂在逐题配对中赢了 8 题、输了 0 题，一题都没输，
额外 prompt token 为 0。留出集复制（另 30 题）得到 0.917，
落在预注册带内（within the pre-registered band，带宽 $0.15$）。

\textbf{oracle 那一轮并没有偏袒上下文臂。} 把金标轮直接交给上下文臂，会被质疑这个对比偏袒上下文：
检索更差会让上下文臂更难看、参数臂更好看。我们把 oracle 换成\textbf{对会话内各轮做相似度检索}，
结果\textbf{测不出任何变化}：精确匹配口径下二者完全相同 —— 0.6B 是 0.517（检索）对
0.517（oracle 选取），1.7B 是 0.533 对 0.533；
冻结地板是 0.017，而把整段历史都塞进去的臂是 0.683；
判官口径下检索臂与 oracle 也在噪声内相当（0.900 对 0.867）。
检索三条轮只能换来 0.550，却要 219 个 prompt token
（单条是 117），所以第二、三条基本是白花的。
这种"打平"是\textbf{短会话子集}的性质；在证据分散于多个会话的多会话题上我们并不指望它保持，
理由见"局限"。

\textbf{同一套流程跑多会话题，以及它测不到什么。} 我们在 30 道\textbf{多会话}题上跑了完全相同的协议，
分数与单会话子集相同（包含率 1.000，地板 0.000，路由 1.000）。
\textbf{这个数字不能被读成"本机制能处理多会话推理"}：本协议里一条记忆就是基准自带的问题配它自己的答案，
所以\textbf{每条都是自足的}，读取路径\textbf{从来不需要把两个会话合起来}，多会话结构根本没到达它。
我们把结果和"它为什么说明不了问题"一起写出来，否则读者自己发现之后就会怀疑协议还藏了什么。

有两条测量教训是我们自己该报告的，它们让表面差距看起来比实际更大。第一，
词面包含判据有偏，它奖励逐字复现，会给同样内容的正确改写判 0。
第二，我们在这个对比上的第一版读数是两种介质分辨不出高下，
当时参数臂的包含率是 0.817。原因是我们自己的生成预算只有 32 token，
把答案截断了，判官于是合理地判它不完整。把所有臂的预算一起提到能写下最长参考答案之后，
参数臂变成 1.000。

在 1.7B 上重跑，参数臂是 1.000，上下文递送是 0.967
到 0.967，成本列不变（0 对最多 1688 个 token），
参数臂仍然一题不输。领先幅度随规模收窄，因为更大的模型更会用上下文：
0.100 的题它能凭先验答对，0.6B 只有 0.017。
成本优势与规模无关，质量优势与规模有关。这是一条两点曲线，我们按它本来的样子写。

\subsection{为什么不干脆把同样的记忆训进权重}
把事实放进上下文是替代记忆库的\textbf{便宜}做法，把它们\textbf{训进权重}是\textbf{昂贵}的那种，
而审稿人一定会问后者。我们跑了三条臂，\textbf{条目-梯度步数完全相同}（30 条 × 12 步；
注意全量微调每一步都更新全部参数，因此它拿到的算力\textbf{更多}）：
全量微调达到 0.133，与记忆库总秩相同的单个适配器是 0.000，
而槽位方案是 1.000。槽位方案对冻结主干也最温和：全位置 KL 为 0.1637 nats，
全量微调是 0.4056、联合适配器是 3.7800；
而且\textbf{只有它能事后忘掉单独一条}：擦掉一个槽后，该槽可证明回到未写过状态
（yes），它的邻居仍然答得出（yes），
而被擦掉的那个问题不再被回答（yes）。
\textbf{必须说明的边界}：三条臂都停在记忆库的预算上，所以这是"\textbf{同等预算}"下的结论；
"微调在更多步数下能到哪"是另一个问题。

\section{负结果}
\label{sec:neg}

\subsection{参数化记忆没有减轻上下文惯性（T5）}
把同一条前提分别放进上下文与权重，然后切换话题。0.6B 上两种做法把前提带进下一话题的程度
相同（1.00 对 1.00），干净场景上的残留也相同（0.42 对
0.42）。在 3 个种子上，这两者\textbf{每一个种子都相等}
（1 对 1）。地板条件是 0.00，也就是哪里都没有前提；
控制臂写了再擦，精确塌回地板。\textbf{换到更大的基座之后两者分开了，而分开的方向让否证更强}：
参数臂仍然继承 1.000，而上下文臂掉到 0.833（地板是 0.000）
—— 也就是说，把前提放进权重\textbf{至少}和放进上下文一样会被带过去，绝不会更少；
擦除控制臂在这个规模上也精确塌回地板。惯性的来源是模型在条件化记忆的内容，
与那些内容是否占着上下文 token 无关。

附带一个值得记下的观察，以及一条诚实说明 —— \textbf{它的方向每个种子都成立，但幅度随种子变化}：
有记忆会让模型更自信。我不知道的比例在无记忆时是 0.167--0.500，
有记忆时降到 0。它不只是被更好地告知，
也更敢直接作答。

\subsection{似然判据量到的是措辞（T3-b）}
这个判据对模型完全答得出的内容判出高惊讶度。逐 token 诊断给出原因：模型从不用裸值作答，
它答 “France's capital is Paris.”，强制续写 \texttt{" Paris"} 要 7.92 nats，
模型自己那句话只要 1.11 nats。两类内容的分布因此几乎重合
（中位数 7.72 对 6.16），实测也差：
漏写 2、误写 6。推广的教训是：只有当探针允许模型
用自己的措辞时，惊讶度才是未知的代理量。

\subsection{两槽崩塌到底是什么（T2-b）}
在动手修理之前，我们先\textbf{不训练任何东西}把它测清楚，共 24 个探针。
写入的更新彼此几乎正交（两两余弦 0.001，最高不超过 0.008），
所以槽位并\textbf{没有在权重方向上互相抢夺}；而对手槽对 logit 的扰动量，达到了它自己那条记忆效果的
很大一部分（0.907），这就是干扰项会致命的原因。
\textbf{秩耗尽也不是原因}：每个槽是按 rank 4 写入的，而在 24 个槽上，
它的能量平均只落在其中 1.97 个方向上（1.88 到 2.05），
也就是说写入\textbf{只用了给它的预算的一半左右}。
但\textbf{根因在别处}。读两个槽\textbf{不等于}把两个单槽效应相加：
实测响应只有该和的 0.847，\textbf{失效是在读取时被制造出来的}。
正确取值 token 的概率，从单选对槽时的 0.861，掉到选中前两槽的 0.695、
以及整个记忆库全开时的 0.032。
写入端的修理必须撤销一个\textbf{只在读取时存在}的相互作用 —— 这正是下面那种修理不奏效的原因。

\subsection{B1 保留项：假设未复现（T8）}
两槽读取的失败与检索无关，正确槽位落在前两名之内的比例是 1.000，
问题出在写入制造的干扰。我们在写入目标里加入保留项（§\ref{sec:method}，
$\lambda_{\mathrm{ret}}$ 取 3.0）：

\begin{table}[h]\centering\small
\begin{tabular}{lccc}\toprule
& 基线 & 加保留项 & 逐探针配对 \\ \midrule
两槽召回 & 0.375 & 0.391 & \textbf{11 帮助 / 10 损害，$p$ 为 1.000} \\
全部槽一起激活 & 0.000 & 0.109 & 7 / 0，$p$ 为 0.016 \\
自身槽位召回 & 0.953 & 0.984 & --- \\
基座损伤（nats） & 1.518 & 0.317 & --- \\ \bottomrule
\end{tabular}
\caption{B1 保留项在 64 个探针上的结果，每条记忆，共 8 条。}
\end{table}

它没有修复两槽召回。在 64 个探针上，配对结果是
11 帮助对 10 损害，精确符号检验的 $p$ 为 1.000，
两槽召回从 0.375 动到 0.391，与噪声没有区别。
更早一次 24 探针的运行看起来像修复（7 对
1，$p$ 为 0.070），把样本翻四倍后没有复现，
所以两次都留在这里。

站得住的是两件事，而且每一次运行都成立。基座损伤降了约五分之四，
从 1.518 到 0.317 nats。保留项真正被评估的那个条件
单向改善：全开档从 0.000 到 0.109，
7 帮助对 0 损害，$p$ 为 0.016。
幅度小，但它在符号检验下是可靠的。总结起来，教写入保护它的前辈，买到了基座损伤的大幅下降
与多槽条件下的小幅改善；我们假设的两槽修复没有证据。

\section{限制}
\label{sec:limits}
\begin{itemize}\setlength\itemsep{2pt}
\item \textbf{规模。}\textbf{每一个头条结论都在 1.7B 上重跑并存活}：组合结论、
写入判据（T3：同样是 5 次写入对 22，
漂移优势从 4.8x \textbf{扩大到} 6.6x）、真实内容诊断、
惯性否证（在更大基座上由"精确相等"变成"\textbf{参数臂更强}"，
1.000 对 0.833，因此结论没有变弱）、
以及遗忘实验（计数完全相同：1/1 对 3/3，
擦除 30/30 个槽逐位验证）。
没有存活的是路由器，它的准确率从 0.925 掉到 0.725，
瓶颈因此从干扰挪到了检索，而没有消失。
仍未测的是 3B 及以上 —— 这台机器装不下。
\item \textbf{基座要付代价。}写入会移动冻结主干。流进 12 条记忆后，
通用问题不再被逐字答对，全位置 KL 约 $1$ nat。参数化记忆是有代价的，
而且代价随写入次数增长。
\item \textbf{B1 那条头条效应没有复现。}24 探针时它像两槽修复
（$p$ 为 0.070），64 探针时配对是 11
对 10（$p$ 为 1.000），也就是什么都没有。我们把两次都留在论文里，
因为效应随样本缩小这件事本身就是结论。
\item \textbf{公开基准只提供范围受限的诊断。}记忆条就是基准自带的问题与答案对，
证据轮由基准自己的标记定位，检索被 oracle 化；写入查询等于读取查询；
判官校准是 12/12（Wilson 区间
0.757–1.000），提示词还改过一次；
判官读得到参考答案，因此残留的偏向接近逐字无法排除。
我们后来也在 \textbf{multi-session 子集}上跑了同一套协议，结果是 1.000，
但\textbf{这个协议测不到多会话结构}（记忆条自足，读取路径不需要跨会话组合，见 §\ref{sec:results}），
所以"单槽无法跨条组合"这条边界在这里仍是\textbf{论证}，还没有变成实测。
\textbf{不过这道边界所依赖的那个事实是测过的}：多会话题中位数带
2.0 条含答案的轮、跨 2.0 个会话，
其中 91\% 跨了不止一个会话；而我们实际评测的单会话子集是
1.0 条轮、1.0 个会话、跨会话比例 0\%。
基于抽取的多会话结果我们\textbf{确实还没有}：抽取调用依赖一个在本批中途被暂停的托管凭据。
\item \textbf{指标粗糙。}基座退化有一部分靠答案是否逐字相同来判断，
它会把 $2+2=4$ 相对冻结态的 2 算作变化；KL 数字更有信息量。
\item \textbf{种子覆盖不均。}T2 跑了五个种子、T3 三个（漂移因此是区间），T4 与 T5 正在
各补到三个；T1 与 T6 是断言型而非统计型，加种子不改变它们。我们没有做\textbf{统一种子矩阵}，
取而代之的是给公开基准诊断做了留出集复制（§\ref{sec:results}）—— 那里才最可能被一次幸运抽样影响。
所以我们只在"数字会动"的地方量化了种子方差，在不会动的地方没有假装量化；
每个实验脚本都接受种子参数。
\end{itemize}

\section{复现方式}
\begin{verbatim}
python -m venv .venv && .venv\Scripts\pip install -e .
python experiments/run_all.py --collect-only   # 从已有报告重建全部数字
python experiments/run_all.py --mode quick     # 冒烟，分钟级
python experiments/run_all.py --mode full      # 从零重跑全部阶段
python paper/build_tables.py                   # 生成英文版表格与宏
python paper/build_paper_zh.py                 # 生成本审阅稿，随后用 xelatex 编译
\end{verbatim}
每个阶段都是独立进程并写自己的 JSON。失败的阶段在汇总里标 \texttt{FAILED}，
其数字渲染为 \texttt{n/a}，绝不用 0 冒充。本文件里的每个数字都来自
\texttt{runs/summary.json}。

\section{结论}
一条秩块槽位的旁路把三件事做得很好，第四件事它\textbf{做不到}。它写入时不打扰邻居，
读取时不花上下文，遗忘可以逐位验证，跨进程重启也成立。它做不到的是\textbf{让多条记忆同时在场}：
把所有槽加起来毫无价值，选两个就已经损失大部分收益，而\textbf{原因已经测出来了}：
两槽的效应不等于单槽效应之和，失效是在\textbf{读时}被制造出来的，我们试过的两种写入端修理都够不到它。
读取时不花上下文，遗忘可以逐位验证，跨进程重启也成立。它不能把模型与上一个话题隔开：
记忆无论存在哪里都有粘性。

我们带走的实用规则有四条。读时组合默认只选一个槽。写入判据应当去问模型；
给似然设阈值会量到措辞。容量应当花在持久且可精确擦除的槽上。我们追过的那一个例外
在相反方向上很有教益：教写入保护记忆库能大幅降低基座损伤、略微改善全槽条件，
但没有买回第二个槽。更小的样本曾暗示它买回了，四倍样本否证了它。
一个随样本量缩小的效应本身就是一个结果。
\end{document}
